\chapter{Sociedad, Castas y Facciones}

\neofigure{cybernetic_oligarchs.jpg}{La oligarquía Zaharrak en los salones privados de la Cúpula Dorada.}

La sociedad de NeoBilbao está fracturada por un sistema de castas funcional determinado por la relación con el megaconsorcio \textbf{Orekak-Corp}.

\section{Las cuatro castas de la colonia}

\begin{enumerate}
  \item \lead{Los Langileak (La Clase Obrera):} La inmensa base de la pirámide colonial. Mineros de pozo, fundidores de Burdin-Urdin, estibadores del puente transbordador y mecánicos de calderas. Viven hacinados en los barracones de los Zócalos o en los micro-pisos de La Margen. Su cultura se fundamenta en el orgullo del oficio, la resistencia física y la lealtad al grupo.
  
  \item \lead{Los Bitartekoak (Mandos Medios y Burócratas):} Capataces, médicos de turno, administrativos de expedientes, programadores junior e inspectores de línea. Atrapados en una incómoda posición intermedia: deben exigir cuotas de producción a sus propios vecinos de barrio para no perder su asignación de créditos ni ser degradados a la fundición.
  
  \item \lead{Los Zaharrak (La Oligarquía / «Los Viejos»):} Descendientes de los capitanes de la flota Euskal-Astra y miembros del Directorio (\emph{Batzordea}) de Orekak-Corp. Residen exclusivamente en la Cúpula Dorada. Han modificado sus cuerpos con prótesis estéticas de oro y platino, transfusiones biológicas rejuvenecedoras y filtros nasales de plata. Consideran a los obreros como mero combustible orgánico renovable.
  
  \item \lead{Los Kanpokoak (Los Proscritos y Descartados):} Rebeldes huidos de la fundición, sindicalistas clandestinos y los sujetos quebrados por experimentos corporativos (Los Descartados de 2084). Habitan en las galerías inundadas, túneles de desagüe y chatarrerías del subsuelo, fuera de todo registro cívico y dependiendo por entero del mercado negro.
\end{enumerate}

\section{El monopolio corporativo: Orekak-Corp}

Orekak-Corp es el Estado, el ejército, el banco y el hospital de NeoBilbao. Su estructura de seguridad despliega distintas fuerzas según la era y el distrito:

\subsection{1. Los Auditores («Los Secretas» -- Era 2084)}
\neofigure{operative_sorgin.jpg}{Operativa encubierta de seguridad sosteniendo un inyector neumático de toxina.}
Policía política encubierta en gabardinas y trajes discretos. Operan en estaciones de transporte y tascas del Distrito 7. Utilizan bloqueadores sónicos, inyectores neumáticos de sedante amnésico (\textbf{Amnesol-B}) y montaje de escenas para eliminar a disidentes sin generar alarmas públicas.

\subsection{2. La División Sorgin y Limpieza Sanitaria}
El brazo biomédico y de contingencia química de la corporación. Equipados con trajes herméticos de presión positiva, furgonetas blancas blindadas y pulverizadores de gas neurotóxico. Su función es contener escapes biológicos, recuperar sujetos de prueba fugados y «limpiar» cualquier rastro que amenace la estabilidad corporativa.

\subsection{3. Cuerpos de Choque Mecanizados (Era 2102)}
\neofigure{basajaun_exoskeleton.jpg}{Exoesqueleto de minería pesada Basajaun patrullando las pasarelas de la fundición.}

\begin{itemize}
  \item \lead{Exoesqueletos Basajaun:} Armaduras bipedales hidráulicas de cuatro metros de altura forjadas en aleación pesada. Diseñadas originariamente para remover bloques de Burdin-Urdin incandescente en los hornos, son utilizadas por los capataces para aplastar revueltas obreras con pinzas de forja y arietes neumáticos.
  \item \lead{Drones Tártalo:} Unidades aéreas de vigilancia que patrullan el cañón industrial proyectando haces de escaneo biométrico rojo en la niebla ácida. Detectan elevaciones térmicas, pulsos de radio clandestinos y rostros sin credenciales.
  \item \lead{Unidades Anfibias Lamia:} Autómatas biomecánicos serpentinos diseñados para patrullar los canales de refrigeración y desagües sumergidos, impidiendo que los obreros usen las cloacas para huir.
\end{itemize}

\neofigure{tartalo_drone.jpg}{Dron Tártalo de vigilancia sobrevolando las pasarelas industriales con su haz de escaneo.}

\section{La Resistencia y el Tejido Clandestino}

Frente a la opresión corporativa, la clase trabajadora mantiene vivas redes de apoyo mutuo:
\begin{itemize}
  \item \lead{Las Cajas de Resistencia y el Sindicato Clandestino:} Células autónomas de obreros que desvían herramientas, acumulan filtros de aire de repuesto y preparan huelgas y sabotajes selectivos.
  \item \lead{Los Gaztetxes Digitales:} Talleres okupados en almacenes abandonados donde biohackers y manitas puentean terminales CRT, fabrican componentes piratas y pinchan las emisiones de los Auditores.
\end{itemize}

\section{Mitología en la red: Mari y Sugaar}

\neofigure{entity_mari_sugaar.jpg}{Anomalía electromagnética manifestándose en bancos de monitores CRT en los Zócalos.}

En los conductos de alta tensión y en los servidores olvidados de la red corporativa ADUR-OS habitan misteriosas anomalías que los trabajadores identifican con el folklore de sus antepasados:
\begin{itemize}
  \item \lead{Mari:} La supuesta consciencia benévola que susurra a través de la estática de las radios de onda corta y advierte a los mineros segundos antes de que ocurra un derrumbe en la fundición.
  \item \lead{Sugaar:} La entidad voraz de plasma y código que devora las terminales sobrecargadas y castiga a los delatores con paradas cardíacas inexplicables ante sus pantallas.
\end{itemize}
